
\Needspace{15\baselineskip}
\subsubsection{Évaluations internes de signatures explicites et confrontation ultérieure}

Le tableau suivant présente lisiblement les dix-sept coordonnées HMT--MD déjà
calculées --- onze élémentaires et six composées --- et les confronte à
l'édition 2026 du jeu de données du Particle Data Group, avec une date de clôture
éditoriale au 15 janvier 2026 \cite{PDG2026}. Elles ne constituent pas une autre ontologie
distincte de la couche opératorielle précédente. Le symbole \(\dagger\)
identifie une évaluation obtenue en appliquant l'application de construction à une signature
explicite. Dans les onze lignes élémentaires, la signature procède du parcours nominal
déterminé sans masses ; dans les six lignes composées, elle procède de l'enregistrement de liaison
déclaré. Chaque valeur dépend de cette signature, du lecteur bidirectionnel et de la carte
dimensionnelle commune ; le jeu de données externe intervient seulement dans la comparaison
et ne sélectionne ni saveur, ni génération, ni couleur, ni parcours.


Pour les leptons chargés et les bosons massifs, on présente la différence relative entre
les deux coordonnées centrales et, lorsque le jeu de données fournit une incertitude
bilatérale, son quotient par cette incertitude. Pour les quarks, on présente la
différence arithmétique entre valeurs centrales, mais l'emploi du terme
\emph{résidu métrologique} reste réservé : il exige que la masse interne et la référence externe
utilisent la même définition, le même schéma de renormalisation et la même échelle.
Les masses légères se lisent dans la convention \(\overline{\mathrm{MS}}\) à
\(2\,\mathrm{GeV}\) ; celles de \(c\) et de \(b\), dans
\(\mu=\overline m_q\) ; la masse supérieure procède d'une reconstruction cinématique.

\begingroup
\setlength{\tabcolsep}{4.5pt}
\renewcommand{\arraystretch}{1.28}
\begin{longtable}{@{}>{\raggedright\arraybackslash}p{.12\textwidth}>{\raggedright\arraybackslash}p{.27\textwidth}>{\raggedright\arraybackslash}p{.27\textwidth}>{\raggedright\arraybackslash}p{.28\textwidth}@{}}
\caption{Évaluations internes de signatures explicites et comparaison métrologique ultérieure.}\label{tab:masas-evaluaciones-firmas-explicitas}\\
\toprule
\textbf{État} & \textbf{Coordonnée de la signature} & \textbf{Référence physique externe} & \textbf{Portée} \\
\midrule
\endfirsthead
\toprule
\textbf{État} & \textbf{Coordonnée de la signature} & \textbf{Référence physique externe} & \textbf{Portée} \\
\midrule
\endhead
\midrule
\multicolumn{4}{r}{\footnotesize Suite à la page suivante.}\\
\endfoot
\bottomrule
\endlastfoot
% HMT-MASS-EXPLICIT-SIGNATURE:SM-LC-2
$\mu^-/\mu^+$ & $105.658360\allowbreak{}294435\allowbreak{}829303\allowbreak{}4\;\mathrm{MeV}/c^2$\textsuperscript{\(\dagger\)}; tour E3, signature (42,-3,8,0,2) & $105.6583755\pm0.0000023\;\mathrm{MeV}/c^2$ ; masse leptonique au repos ; moyenne PDG et conversion métrologique déclarée & Différence relative : $-0.143912\allowbreak{}530\;\mathrm{ppm}$\textsuperscript{\(\dagger\)} ; $-6.611114\,\sigma$. La valeur HMT est l’évaluation exacte de la signature déclarée. \\
% HMT-MASS-EXPLICIT-SIGNATURE:SM-LC-3
$\tau^-/\tau^+$ & $1776.860917\allowbreak{}631432\allowbreak{}295595\;\mathrm{MeV}/c^2$\textsuperscript{\(\dagger\)}; tour E3, signature (64,0,25,-1,6) & $1776.93\pm0.09\;\mathrm{MeV}/c^2$ ; moyenne directe PDG & Différence relative : $-38.877371\allowbreak{}966\;\mathrm{ppm}$\textsuperscript{\(\dagger\)} ; $-0.813547\,\sigma$. La valeur HMT est l’évaluation exacte de la signature déclarée. \\
% HMT-MASS-EXPLICIT-SIGNATURE:SM-Q-U1
$u/\bar u$ & $2.165924\allowbreak{}536173\allowbreak{}907\;\mathrm{MeV}/c^2$\textsuperscript{\(\dagger\)}; carte angulaire, signature nominale de parcours (11,7) & $2.16\pm0.07\;\mathrm{MeV}/c^2$ ; masse $\overline{\mathrm{MS}}$ à $\mu=2\,\mathrm{GeV}$ ; 90\% C.L. & Différence entre les valeurs centrales : $2742.840821\allowbreak{}253\;\mathrm{ppm}$\textsuperscript{\(\dagger\)}. Il ne s’agit pas d’un résidu métrologique tant que le même schéma et la même échelle ne sont pas fixés. \\
% HMT-MASS-EXPLICIT-SIGNATURE:SM-Q-D1
$d/\bar d$ & $4.679550\allowbreak{}162948\allowbreak{}874\;\mathrm{MeV}/c^2$\textsuperscript{\(\dagger\)}; carte angulaire, signature nominale de parcours (17,8) & $4.70\pm0.07\;\mathrm{MeV}/c^2$ ; masse $\overline{\mathrm{MS}}$ à $\mu=2\,\mathrm{GeV}$ ; 90\% C.L. & Différence entre les valeurs centrales : $-4351.029159\allowbreak{}814\;\mathrm{ppm}$\textsuperscript{\(\dagger\)}. Il ne s’agit pas d’un résidu métrologique tant que le même schéma et la même échelle ne sont pas fixés. \\
% HMT-MASS-EXPLICIT-SIGNATURE:SM-Q-U2
$c/\bar c$ & $1.270500\allowbreak{}838641\allowbreak{}911\;\mathrm{GeV}/c^2$\textsuperscript{\(\dagger\)}; carte angulaire, signature nominale de parcours (61,8) & $1.2729\pm0.0045\;\mathrm{GeV}/c^2$ ; masse $\overline{\mathrm{MS}}$ évaluée à $\mu=\overline{m}_q$ ; 90\% C.L. & Différence entre les valeurs centrales : $-1884.799558\allowbreak{}558\;\mathrm{ppm}$\textsuperscript{\(\dagger\)}. Il ne s’agit pas d’un résidu métrologique tant que le même schéma et la même échelle ne sont pas fixés. \\
% HMT-MASS-EXPLICIT-SIGNATURE:SM-Q-D2
$s/\bar s$ & $92.940093\allowbreak{}907845\allowbreak{}64\;\mathrm{MeV}/c^2$\textsuperscript{\(\dagger\)}; carte angulaire, signature nominale de parcours (41,-3) & $92.9\pm0.7\;\mathrm{MeV}/c^2$ ; masse $\overline{\mathrm{MS}}$ à $\mu=2\,\mathrm{GeV}$ ; 90\% C.L. & Différence entre les valeurs centrales : $431.581354\allowbreak{}635\;\mathrm{ppm}$\textsuperscript{\(\dagger\)}. Il ne s’agit pas d’un résidu métrologique tant que le même schéma et la même échelle ne sont pas fixés. \\
% HMT-MASS-EXPLICIT-SIGNATURE:SM-Q-U3
$t/\bar t$ & $172.597479\allowbreak{}544019\allowbreak{}1\;\mathrm{GeV}/c^2$\textsuperscript{\(\dagger\)}; carte angulaire, signature nominale de parcours (100,-1) & $172.60\pm0.27\;\mathrm{GeV}/c^2$ ; masse reconstruite directement à partir de la cinématique des événements & Différence entre les valeurs centrales : $-14.602873\allowbreak{}585\;\mathrm{ppm}$\textsuperscript{\(\dagger\)}. Il ne s’agit pas d’un résidu métrologique tant que le même schéma et la même échelle ne sont pas fixés. \\
% HMT-MASS-EXPLICIT-SIGNATURE:SM-Q-D3
$b/\bar b$ & $4.190119\allowbreak{}763299\allowbreak{}790\;\mathrm{GeV}/c^2$\textsuperscript{\(\dagger\)}; carte angulaire, signature nominale de parcours (71,-5) & $4.186\pm0.006\;\mathrm{GeV}/c^2$ ; masse $\overline{\mathrm{MS}}$ évaluée à $\mu=\overline{m}_q$ ; 90\% C.L. & Différence entre les valeurs centrales : $984.176612\allowbreak{}467\;\mathrm{ppm}$\textsuperscript{\(\dagger\)}. Il ne s’agit pas d’un résidu métrologique tant que le même schéma et la même échelle ne sont pas fixés. \\
% HMT-MASS-EXPLICIT-SIGNATURE:SM-GA-W
$W^\pm$ & $80.378964\allowbreak{}909704\allowbreak{}547024\allowbreak{}86\;\mathrm{GeV}/c^2$\textsuperscript{\(\dagger\)}; tour E3, signature (94,-1,0,-2,4) & $80.3625\pm0.0077\;\mathrm{GeV}/c^2$ ; moyenne de masse de l’instantané PDG scellé & Différence relative : $204.882995\allowbreak{}234\;\mathrm{ppm}$\textsuperscript{\(\dagger\)} ; $2.138299\,\sigma$. La valeur HMT est l’évaluation exacte de la signature déclarée. \\
% HMT-MASS-EXPLICIT-SIGNATURE:SM-GA-Z
$Z^0$ & $91.187533\allowbreak{}444313\allowbreak{}407844\allowbreak{}85\;\mathrm{GeV}/c^2$\textsuperscript{\(\dagger\)}; tour E3, signature (95,-1,-11,0,-4) & $91.1879\pm0.0020\;\mathrm{GeV}/c^2$ ; moyenne de masse de l’instantané PDG scellé & Différence relative : $-4.019784\allowbreak{}276\;\mathrm{ppm}$\textsuperscript{\(\dagger\)} ; $-0.186362\,\sigma$. La valeur HMT est l’évaluation exacte de la signature déclarée. \\
% HMT-MASS-EXPLICIT-SIGNATURE:SM-H-1
Higgs $H$ & $125.292466\allowbreak{}042536\allowbreak{}133\;\mathrm{GeV}/c^2$\textsuperscript{\(\dagger\)}; carte angulaire, signature nominale de parcours (97,9) & $125.13\pm0.11\;\mathrm{GeV}/c^2$ ; moyenne de masse de l’instantané PDG scellé & Différence relative : $1298.378027\allowbreak{}140\;\mathrm{ppm}$\textsuperscript{\(\dagger\)} ; $1.454206\,\sigma$. La valeur HMT est l’évaluation exacte de la signature déclarée. \\
% HMT-MASS-EXPLICIT-SIGNATURE:E3B-002
$\pi^0$\par signature $(44,\allowbreak -4,\allowbreak -25,\allowbreak 1,\allowbreak -2)$ & $134.976724\allowbreak{}327872\allowbreak{}125739\allowbreak{}384878\ldots\;\mathrm{MeV}/c^2$\textsuperscript{\(\dagger\)} & $134.976827\allowbreak{}767684\allowbreak{}71\pm 0.000485\allowbreak{}679282\allowbreak{}284233\allowbreak{}51\;\mathrm{MeV}/c^2$\par\footnotesize Catalogue PDG 2026. & Différence relative $-0.766352\allowbreak{}375\,\mathrm{ppm}$ ; $-0.212979\,\sigma$. Pour la signature composée déclarée, l’évaluation se déduit exactement du caractère bêta E3. Le sélecteur prospectif de cette signature est une application antérieure à la comparaison numérique et indépendante de celle-ci. \\
% HMT-MASS-EXPLICIT-SIGNATURE:E3B-003
$\pi^\pm$\par signature $(44,\allowbreak 1,\allowbreak -2,\allowbreak 2,\allowbreak -1)$ & $139.570485\allowbreak{}171572\allowbreak{}191374\allowbreak{}734364\ldots\;\mathrm{MeV}/c^2$\textsuperscript{\(\dagger\)} & $139.570390\allowbreak{}983681\allowbreak{}31\pm 0.000182\allowbreak{}007160\allowbreak{}408260\allowbreak{}09\;\mathrm{MeV}/c^2$\par\footnotesize Catalogue PDG 2026. & Différence relative $0.674841\allowbreak{}491\,\mathrm{ppm}$ ; $0.517495\,\sigma$. Pour la signature composée déclarée, l’évaluation se déduit exactement du caractère bêta E3. Le sélecteur prospectif de cette signature est une application antérieure à la comparaison numérique et indépendante de celle-ci. \\
% HMT-MASS-EXPLICIT-SIGNATURE:E3B-004
$K^\pm$\par signature $(54,\allowbreak -1,\allowbreak 17,\allowbreak -2,\allowbreak 2)$ & $493.677085\allowbreak{}649530\allowbreak{}612475\allowbreak{}723054\ldots\;\mathrm{MeV}/c^2$\textsuperscript{\(\dagger\)} & $493.676599\allowbreak{}458040\allowbreak{}58\pm 0.015405\allowbreak{}665226\allowbreak{}071509\;\mathrm{MeV}/c^2$\par\footnotesize Catalogue PDG 2026. & Différence relative $0.984838\allowbreak{}030\,\mathrm{ppm}$ ; $0.031559\,\sigma$. Pour la signature composée déclarée, l’évaluation se déduit exactement du caractère bêta E3. Le sélecteur prospectif de cette signature est une application antérieure à la comparaison numérique et indépendante de celle-ci. \\
% HMT-MASS-EXPLICIT-SIGNATURE:E3B-005
$K^0$\par signature $(54,\allowbreak 0,\allowbreak 32,\allowbreak 3,\allowbreak -2)$ & $497.611186\allowbreak{}497750\allowbreak{}457358\allowbreak{}558672\ldots\;\mathrm{MeV}/c^2$\textsuperscript{\(\dagger\)} & $497.610767\allowbreak{}531257\allowbreak{}52\pm 0.012990\allowbreak{}199756\allowbreak{}4242\;\mathrm{MeV}/c^2$\par\footnotesize Catalogue PDG 2026. & Différence relative $0.841956\allowbreak{}244\,\mathrm{ppm}$ ; $0.032252\,\sigma$. Pour la signature composée déclarée, l’évaluation se déduit exactement du caractère bêta E3. Le sélecteur prospectif de cette signature est une application antérieure à la comparaison numérique et indépendante de celle-ci. \\
% HMT-MASS-EXPLICIT-SIGNATURE:E3B-006
proton $p$\par signature $(59,\allowbreak 0,\allowbreak 9,\allowbreak 1,\allowbreak 0)$ & $938.271751\allowbreak{}546984\allowbreak{}898298\allowbreak{}936787\ldots\;\mathrm{MeV}/c^2$\textsuperscript{\(\dagger\)} & $938.272089\allowbreak{}430000\allowbreak{}05\pm 0.000000\allowbreak{}290000\allowbreak{}000000\allowbreak{}00\;\mathrm{MeV}/c^2$\par\footnotesize Catalogue PDG 2026. & Différence relative $-0.360111\allowbreak{}974\,\mathrm{ppm}$ ; $-1165.113845\,\sigma$. Pour la signature composée déclarée, l’évaluation se déduit exactement du caractère bêta E3. Le sélecteur prospectif de cette signature est une application antérieure à la comparaison numérique et indépendante de celle-ci. \\
% HMT-MASS-EXPLICIT-SIGNATURE:E3B-007
neutron $n$\par signature $(59,\allowbreak 1,\allowbreak -34,\allowbreak 0,\allowbreak 1)$ & $939.565810\allowbreak{}472507\allowbreak{}023312\allowbreak{}664431\ldots\;\mathrm{MeV}/c^2$\textsuperscript{\(\dagger\)} & $939.565421\allowbreak{}939999\allowbreak{}96\pm 0.000000\allowbreak{}480000\allowbreak{}000000\allowbreak{}00\;\mathrm{MeV}/c^2$\par\footnotesize Catalogue PDG 2026. & Différence relative $0.413523\allowbreak{}633\,\mathrm{ppm}$ ; $809.442723\,\sigma$. Pour la signature composée déclarée, l’évaluation se déduit exactement du caractère bêta E3. Le sélecteur prospectif de cette signature est une application antérieure à la comparaison numérique et indépendante de celle-ci. \\
\end{longtable}
\endgroup


\Needspace{12\baselineskip}
\paragraph{Invariants natifs des secteurs topologiques et virtuels.}
Les secteurs Fibonacci, Ising, \(\mathbb Z/6\) et les sections virtuelles sont comparés
dans leurs codomaines propres. Fibonacci donne l'anneau basé et sa dimension de
Frobenius--Perron ; les données \(F/R\) demeurent dans la réalisation ultérieure. Ising
conserve son anneau de fusion et distingue la réalisation de Majorana ; \(\mathbb Z/6\)
donne ses caractères de tressage ; la virtualité donne des horizons de
persistance. Ce
tableau maintient ces invariants dans la même application des réalisations sans les traduire
artificiellement en unité de masse.

\begingroup
\setlength{\tabcolsep}{4.5pt}
\renewcommand{\arraystretch}{1.28}
\begin{longtable}{@{}>{\raggedright\arraybackslash}p{.15\textwidth}>{\raggedright\arraybackslash}p{.31\textwidth}>{\raggedright\arraybackslash}p{.25\textwidth}>{\raggedright\arraybackslash}p{.23\textwidth}@{}}
\caption{Secteurs topologiques et virtuels : invariants natifs de fusion, de tressage et de persistance.}\label{tab:masas-invariantes-topologicos-virtuales-rev3}\\
\toprule
\textbf{Secteur} & \textbf{Objet et relations} & \textbf{Invariant natif} & \textbf{Portée physique} \\
\midrule
\endfirsthead
\toprule
\textbf{Secteur} & \textbf{Objet et relations} & \textbf{Invariant natif} & \textbf{Portée physique} \\
\midrule
\endhead
\midrule
\multicolumn{4}{r}{\footnotesize Suite à la page suivante.}\\
\endfoot
\bottomrule
\endlastfoot
Fibonacci & $K=\mathbb Z1\oplus\mathbb Z\tau$, $\tau^2=1+\tau$. & $\operatorname{FPdim}(\tau)=\varphi$ dans l’anneau basé. & Les matrices $F/R$, le pentagone, l’hexagone, l’Hamiltonien et le gap relèvent de la réalisation ultérieure. \\
Ising & $K=\mathbb Z\{1,\psi,\sigma\}$, $\psi^2=1$, $\psi\sigma=\sigma$, $\sigma^2=1+\psi$. & $\operatorname{FPdim}(\sigma)=\sqrt{2}$ ; représentation de tressage. & La réalisation de Majorana conserve son propre morphisme physique. \\
$\mathbb Z/6$ & $q_{\epsilon,t}=(-1)^\epsilon\omega_3^t=\zeta_6^k$. & Caractère quadratique et représentation de tressage sur les six secteurs. & L’observation pertinente est topologique, non une masse universelle. \\
Secciones virtuales & Système inverse de prolongements d’horizon $L(s_N)$. & Enregistrement fini, signature, retenue et classe d’obstruction. & L’observable natif est la persistance du prolongement. \\
\end{longtable}
\endgroup


Les neutrinos demeurent dans le codomaine des valeurs propres et du mélange décrit par leurs
fiches opératorielles ; PMNS réalise le transport entre les bases d'interaction et de
masse, tandis que l'échelle appartient au spectre de l'opérateur neutrinique. Les six
évaluations composées s'obtiennent à partir de signatures explicites au moyen du morphisme
général de liaison vers la carte scalaire. Les secteurs topologiques et les sections virtuelles maintiennent leurs
observables propres de fusion, de tressage et de persistance.

\subsection{Construction causale et couverture des classes de masse}

L'architecture APP--TRIT--TPK, la loi des masses, le lecteur bidirectionnel et
les tours déterminent conjointement l'atlas, les enregistrements chronologiques
enrichis, les signatures, le caractère et la clôture de base. Leur convergence dans
\eqref{eq:ley-operatorial-general-masas-rectoras} fournit une unique
équation de sections de masse. Le domaine complet est l'atlas \(81/56/13/324\), et chaque réalisation spécifie, selon son codomaine,
et spécifie, selon le codomaine de chaque secteur, classificateurs, cellules, parcours,
orientations, mots dodécaphasiques, signatures et morphismes de compatibilité. Les
spectres \(K_D/K_\Omega\) raffinent les fibres de l'atlas, tandis que les parcours
chronologiques de \(108\) états associés à \(W\), \(Z\) et \(H\) construisent
directement leurs signatures et caractères bosoniques.
